\documentclass[11pt,a4paper]{article}
\usepackage[T1]{fontenc}
\usepackage[utf8]{inputenc}
\usepackage[margin=20mm]{geometry}
\usepackage{textcomp,xurl,microtype}
\usepackage[hidelinks]{hyperref}
\setlength{\parindent}{0pt}
\setlength{\parskip}{5pt}
\setlength{\emergencystretch}{2em}
\begin{document}
23 October 2026\par
The Editors, \textit{Briefings in Bioinformatics}\par
Article type: Problem solving protocol\par
Manuscript: \textit{Separating reconstruction, generation and antibacterial activity in antimicrobial peptide design}\par
Authors: Hongsheng Wang, Fei Wu and Feng Lin

Dear Editors,

We submit this manuscript as a Problem solving protocol for evaluating antimicrobial peptide design. It addresses a practical interpretation problem: sequence reconstruction, de novo generation and experimental antibacterial activity provide distinct evidence, and an improvement in one endpoint does not automatically establish improvement in another.

A matched reconstruction-objective comparison increases development residue accuracy from 12.55\% to 70.90\% with categorical rather than squared-error training. The categorical model also improves six sampled-sequence distances across terminal conditions and reference groups. Direct generation evaluation adds an informative control: a locally trained terminal-conditioned LSTM is closest on all 24 distance comparisons. The protocol therefore makes the contribution of the objective change clear while showing why sampled outputs need their own evaluation.

In a separate panel of 27 experimentally tested peptides, six inhibit \textit{S. aureus} ATCC 29213 at MICs $\leq$31.25 \textmu{}g/mL, and seven have MICs below 200 \textmu{}g/mL against at least one of four reference strains. These counts are stable across the original MIC summary and a replicate-based sensitivity analysis. Complete peptide-level MICs and optical-density measurements accompany the manuscript.

The study complements evaluation-focused work in \textit{Briefings in Bioinformatics} by providing an explicit mapping from the scientific question to its input, comparator and measured endpoint. AMPGen is discussed as a published example of staged generation, selection and biological assessment. Hongsheng Wang and Feng Lin are also authors of AMPGen; the present computational comparison and experimental panel are separate from that publication.

Supplementary Methods, complete tables, model code, checkpoints and executable analysis scripts are available at \href{https://github.com/wanghongsheng01/bib-amp-review-materials-20261023/tree/v5-2026-10-23}{the versioned GitHub repository}. The model comparison is descriptive, using one fitted seed and a development partition used during model development; these conditions are specified in the manuscript.

The authors wrote the manuscript. OpenAI Codex assisted with source and reference checking, code inspection, recalculation of saved outputs, deterministic plotting, factual correction and language editing. The manuscript and supplement disclose this assistance. The author contributions, institutional affiliations and competing-interest statement are included. Fei Wu is named as guarantor. The work involves no human or animal subjects.

Sincerely,\par
\textbf{Hongsheng Wang}\par
Corresponding author\par
College of Computer Science and Technology, Zhejiang University\par
Artificial Intelligence Department, Loctek Ergonomic Technology Corp.\par
\href{mailto:wanghongsheng@zju.edu.cn}{wanghongsheng@zju.edu.cn}
\end{document}
